\subsection{LINEAR EQUATIONS IN VECTOR SPACES}

PROPOSITION 12. Let E, F be two vector spaces over a field K and u:E → F a linear mapping. For the linear equation

\[
u (x) = y _ {0}\tag{20}
\]

to have at least one solution \(x \in \mathbf{E}\) , it is necessary and sufficient that \(y_0\) be orthogonal to the kernel of the transpose mapping \(^t u\) .

The orthogonal of \(u(\mathbf{E})\) in \(\mathbf{F}^*\) is \({}^t u^{-1}(0)\) (§ 2, no. 5, Corollary to Proposition 8) and the orthogonal of \({}^t u^{-1}(0)\) in \(\mathbf{F}\) is therefore \(u(\mathbf{E})\) (no. 5, Theorem 7 (ii)).

We shall obtain a more convenient criterion for systems of scalar linear equations

\[
\langle x, x _ {\iota} ^ {*} \rangle = \eta_ {\iota} \quad (\iota \in I)\tag{21}
\]

where the unknown x takes its values in a vector space E over a field K, the \(x_{\iota}^{*}\) are linear forms on E and the right hand sides \(\eta_{\iota}\) elements of K.

If we consider a basis \((a_{\lambda})_{\lambda\in L}\) of E, the system (21) is equivalent to the system of equations

\[
\sum_ {\lambda \in L} \xi_ {\lambda} \left\langle a _ {\lambda}, x _ {\iota} ^ {*} \right\rangle = \eta_ {\iota} \quad (\iota \in I)\tag{22}
\]

with \(x = \sum_{\lambda \in L} \xi_{\lambda} a_{\lambda}\) , the solutions of (22) being of necessity families \((\xi_{\lambda})\) of elements of \(K\) of finite support.

DEFINITION 4. The dimension of the subspace of \(\mathbf{E}^*\) generated by the family \((x_{\iota}^{*})\) is called the rank of the system (21).

PROPOSITION 13. For the system (21) to be of finite rank \(r\) , it is necessary and sufficient that the linear mapping \(u: x \mapsto (\langle x, x_{\iota}^* \rangle)\) of \(\mathbf{E}\) into \(\mathbf{K}_s^I\) be of rank \(r\) .

If \(F'\) is the subspace of \(E^{*}\) generated by the \(x_{\iota}^{*}\) , the kernel of u is the orthogonal F of \(F'\) in E; if \(F'\) is of dimension r, F is of codimension r and conversely (no. 5, Theorem 7) and \(\operatorname{rg}(u) = \operatorname{codim}_{\mathbf{E}} \mathbf{F}\) (no. 4, formula (11)).

THEOREM 8. Let

\[
\langle x, x _ {\iota} ^ {*} \rangle = \eta_ {\iota} \quad (\iota \in I)\tag{21}
\]

be a system of scalar linear equations on a vector space E over a field K. For this system to have at least one solution, it is necessary that, for every family \((\rho_{\iota})\) of scalars of finite support such that \(\sum_{\iota} x_{\iota}^{*}\rho_{\iota}=0\) , \(\sum_{\iota}\eta_{\iota}\rho_{\iota}=0\) . If the rank of the system (21) is finite, this condition is also sufficient.

The condition is obviously necessary. It says that, if \(F'\) is the subspace of \(E^{*}\) generated by the family \((x_{i}^{*})\) , there exists a linear mapping \(f:F'\to K_{d}\) such that \(f(x_{\iota}^{*})=\eta_{\iota}\) for all \(\iota\in I\) . If \(F'\) is of finite dimension r, \(F'\) is the orthogonal of a subspace F of E of codimension r (no. 5, Theorem 7) and \(F'\) is identified with the dual of E/F (§2, no. 6, Corollary to Proposition 9); f is therefore an element of the bidual \((\mathrm{E}/\mathrm{F})^{**}\) . As E/F is finite-dimensional there exists one and only one element \(y\in E/F\) such that \(f(x^{*})=\langle y,x^{*}\rangle\) for all \(x^{*}\in F'\) (no. 5, Theorem 6). The solutions of (21) are then the \(x\in E\) whose canonical image in E/F is y.

Remark. When the rank of the system (21) is infinite, the condition of Theorem 8 is no longer sufficient. For example, suppose that the \(x_{\iota}^{*}\) are the coordinate forms on the space \(\mathbf{E} = \mathbf{K}_{s}^{(\mathrm{I})}\) , I being infinite (§ 2, no. 6); as the \(x_{\iota}^{*}\) are linearly independent, the condition of Theorem 8 holds for every family ( \(\eta_{\iota}\) ) but the system (21) only has solutions if the family ( \(\eta_{\iota}\) ) has finite support.

A system (21) is always of finite rank if it has only a finite number of equations and its rank is then at most equal to the number of equations (no. 2, Proposition 1). Similarly, if E is of finite dimension n (which for a system (22) corresponds to the case where there is only a finite number n of unknowns), its dual \(E^{*}\) is of dimension n and hence the rank of system (21) is at most equal to n (no. 3, Corollary 4 to Proposition 4). From this we deduce:

COROLLARY 1. A system of scalar linear equations in a vector space, consisting of a finite number of equations whose left hand sides are linearly independent forms, always admits solutions.

COROLLARY 2. For a homogeneous linear system (22) of equations in n unknowns with coefficients in a field K to admit non-trivial solutions consisting of elements of K, it is necessary and sufficient that its rank be \textless n.

This will always be the case if the number of equations is finite and \textless n.

COROLLARY 3. For a linear system (22) with coefficients and right hand sides in a field K, consisting of n equations in n unknowns, to have one and only one solution consisting of elements of K, it is necessary and sufficient that the associated homogeneous system have no non-trivial solution (or, what amounts to the same, that the left hand sides of this system be linearly independent forms).

